\documentclass[11pt,a4paper]{article}
\usepackage[T1]{fontenc}
\usepackage[utf8]{inputenc}
\usepackage[margin=2.2cm]{geometry}
\usepackage{amsmath,amssymb}
\usepackage{xcolor}
\usepackage{titlesec}

\titleformat{\section}{\normalfont\large\bfseries\color{blue!50!black}}
            {\thesection.}{0.5em}{}
\titlespacing*{\section}{0pt}{1.0em}{0.3em}

\setlength{\parskip}{0.35em}
\setlength{\parindent}{0pt}

\title{\vspace{-2em}\textbf{0--k BFS -- Speaker Script}\\
\large Companion text to \texttt{presentation\_en.pdf}}
\author{Artemov Makar, Mark Shkut \\ Skoltech, Algorithms 2026}
\date{May 21, 2026}

\begin{document}
\maketitle
\vspace{-1em}

The text below is the literal script -- what to say on each slide. Each
section is one slide; read the paragraph aloud and click forward when you
hit the dashed line.

\section{Title}
``Good morning. I'm Makar, this is Mark. We've built a C++17
implementation of the 0--k BFS family of shortest-path algorithms and
benchmarked it against Dijkstra. The headline number is that on a graph
with nine hundred million vertices, our algorithm runs in two minutes
where Dijkstra takes ten. I'll explain the algorithm, show it running,
and walk through the measurements. About fifteen minutes plus questions.''
\hrulefill

\section{The problem}
``The problem is single-source shortest path on a weighted graph. The
twist is that the weights are non-negative integers bounded by a small
known constant $k$. This restriction is not artificial -- it comes up
constantly in grid path planning, in layered graphs, in 2-SAT-style
reductions, in transport networks with a few toll levels. Our concrete
target was to push the input size as close to one billion vertices as
the hardware allows, on a single workstation.''
\hrulefill

\section{Baseline: Dijkstra}
``The classical answer is Dijkstra. We use the textbook implementation
with a \texttt{std::priority\_queue} of (distance, vertex) pairs and lazy
decrease-key -- whenever we improve a distance we push a fresh entry, and
we skip stale entries when we pop them. Complexity $O((V+E)\log V)$.
This works for any non-negative real weights. The $\log V$ factor sits
inside every successful relaxation -- it's the cost of the heap. The
question is, can we get rid of it when the weights are small integers?''
\hrulefill

\section{The trick: 0--1 BFS}
``The simplest case is weights zero or one. We replace the priority
queue with a deque. An edge of weight zero pushes to the front, an edge
of weight one pushes to the back. The deque is always sorted by distance
and contains at most two distinct values at any time -- the current
layer at the front, the next layer at the back. Linear time, linear
memory. One subtle point: unlike plain BFS, a vertex can be discovered
several times before its distance is final, so the relaxation check
must be `dist[v] plus weight is less than dist[u]', not just
`visited or not'.''
\hrulefill

\section{Generalize: 0--k BFS}
``Now generalize to weights up to $k$. The observation is that a
relaxation can move a vertex at most $k$ steps forward in distance
order. So at any moment, only $k+1$ distance layers are alive. We
allocate $k+1$ FIFO queues and index them by distance modulo $k+1$ --
that's the circular buffer of buckets. A pointer \texttt{cur} walks
through them. Relaxation pushes the target into the bucket for the new
distance. When the active bucket empties, \texttt{cur} advances. Stale
entries from earlier relaxations sit in their old buckets and get
skipped on pop. Complexity $O(kV+E)$, memory $O(k+E)$ for the queues
plus $O(V)$ for the distance array.''
\hrulefill

\section{What it looks like}
``Here are both algorithms running on the same grid -- two hundred by
two hundred, $k=8$. They produce the same distance field -- the
wavefront spreads radially from the centre, shaped by random terrain
weights. The titles show the actual measured wall-clock times. The
left panel, 0--k BFS, finishes in four milliseconds. The right panel,
Dijkstra, takes thirteen point three milliseconds -- about three and a
third times slower. Same correctness, much less time. The animation
itself is in the repository; this is one frame from it.''
\hrulefill

\section{Under the hood}
``This is the pedagogical heart. On a tiny grid -- ten by ten,
$k=3$ -- we instrumented the algorithm to dump every event. On the
right you see the four real queues, $Q[0]$ through $Q[3]$, with their
current contents. The numbers inside the cells are vertex ids. On the
left is the grid; each cell's colour shows which bucket currently owns
that vertex. The red `cur' arrow marks the active bucket. Watch how a
relaxation lands a vertex in the bucket chosen by the new distance
modulo four, and how the active bucket drains in FIFO order before
\texttt{cur} advances. That's the whole secret of the algorithm.''
\hrulefill

\section{Designing for V $\to 10^9$}
``Now the engineering. To hit one billion vertices we cannot store edges
explicitly -- a CSR with four edges per vertex and 32-bit ids would
already cost twenty gigabytes. So we built two backends. The general
\texttt{GraphCSR} for ordinary graphs, and an implicit \texttt{GridGraph}
that computes the four neighbours of each cell on the fly with no edge
storage. Second, we made the distance type \texttt{uint32\_t} instead of
\texttt{uint64\_t} -- that halves the dominant memory cost. The
trade-off is only valid as long as $k$ times $V$ fits in a 32-bit
integer, which is fine for everything we care about and we have a
64-bit fallback.''
\hrulefill

\section{Results}
``Here are the numbers, all grid, single thread. At a million vertices
we already see a 4.6$\times$ speed-up. At ten million, 5.6$\times$. At
one hundred million with $k=1$, the deque-based 0--1 BFS reaches an
eleven times speed-up -- this is the peak. The next row is interesting:
at half a billion vertices both algorithms become memory-bandwidth bound
and the gap collapses to about one. The bottom row is the headline -- on
nine hundred million vertices, 0--1 BFS does the work in one hundred
twenty-one seconds where Dijkstra takes six hundred twenty. Every single
distance checksum matches across algorithms -- nineteen distinct graph
instances in the validation sweep, zero mismatches.''
\hrulefill

\section{Scaling}
``Same data, log-log plot. Both axes logarithmic. Dijkstra is the top
line, the BFS variants run parallel underneath. The constant vertical
gap on a log scale is a constant multiplicative speed-up. Towards the
right side, at $10^8$ and beyond, the lines start to converge -- we are
entering the memory-bound regime. The lowest line throughout is the
deque-based 0--1 BFS, which stays cleanly ahead.''
\hrulefill

\section{Limits}
``Honest about the boundaries. First: this is a constant-factor win, not
asymptotic. Once $k$ grows comparable to $\log V$ Dijkstra is at least
as good. Second: above half a billion vertices everything is memory-bound,
and the queue overhead actually starts to hurt. Third: maintaining
$k+1$ separate FIFO allocators is more expensive for the cache than a
single deque -- that's why our 0--1 BFS at $V \approx 10^9$ still wins
but a generic 0--k BFS at the same scale would not. Fourth: this isn't
a replacement for Dijkstra in general -- for real-valued weights or an
unknown weight range, use Dijkstra.''
\hrulefill

\section{Next steps}
``Four planned items. First, a proper indexed $d$-ary heap with real
decrease-key for the baseline -- to be a fairer reference. Second, a
chunked ring-buffer FIFO for 0--k BFS to close the cache gap against
0--1 BFS. Third, leave the synthetic-graph comfort zone and load a real
road network -- DIMACS USA with weights quantised to $k=8$. Fourth, a
stretch goal: a $\Delta$-stepping variant for parallel context.''
\hrulefill

\section{Thank you}
``All of this -- code, the LaTeX report, the animations, reproducible
build scripts -- lives in the GitHub repository linked here. Happy to
take questions.''

\textbf{Likely questions:}
\begin{itemize}
\item \emph{Why is 0--1 BFS faster than 0--k BFS at $k=1$?} -- One
\texttt{std::deque} versus two \texttt{std::queue}s: allocator and cache
overhead.
\item \emph{How did you verify correctness on a billion-vertex run?} --
Distance-sum checksums match Dijkstra exactly; the sum pins down the
array without per-vertex comparison.
\item \emph{What about parallelism?} -- Single thread throughout Stage~1;
$\Delta$-stepping is planned.
\end{itemize}

\end{document}
